\documentclass{article}
\usepackage{etoolbox}
\usepackage{graphicx}

\ifdef{\myoutput}{}{\def\myoutput{concept}}
\ifdef{\myversion}{}{\def\myversion{1}}

\usepackage[output=\myoutput
,numberofversions=4
,version=\myversion
,seed=4
,randomizequestions=true
,randomizeanswers=true
,writeRfile=true
]{mcexam}

\usepackage[margin=0.95in]{geometry}
\usepackage{multicol}

\usepackage{fancyhdr,lastpage}
\pagestyle{fancy}
\fancyhf{}
\renewcommand{\headrulewidth}{1pt} 
\renewcommand{\footrulewidth}{1pt}
\lfoot{\mctheversion}
\rfoot{Page \thepage\ of \pageref{LastPage}}
\lhead{Quiz 4.2 \& 4.3 \mctheversion \\ MAT 2000: Introduction to Statistics}
\rhead{Name Print): \hspace{2.5 in}}

\usepackage{calc}
\renewenvironment{setmcquestion}
{\begin{minipage}[t]{\linewidth-\labelwidth}}
	{\end{minipage}\par}
	
\renewenvironment{setmcquestioninstruction}
{\begin{minipage}{\textwidth}}
	{\end{minipage}}
	
\renewenvironment{setmcanswers}{}{}
\setlist[setmcquestions]{label=\mcquestionlabelfmt{*}.
	,ref=\mcquestionlabelfmt{*}
	,itemsep=0.5\baselineskip
	,topsep=1\baselineskip
}

\renewcommand\mcversionlabelfmt[1]{\arabic{#1}}

%copy-pasting from MS Word, use regex to find [\*]*[a-d]\. and replace with \answer 

\raggedbottom

\begin{document}
	Circle the correct answer for each question.  If you're unsure, you can write numbers 1, 2, 3, etc. next to the choices to get partial credit.\\[.5cm]
	
	A 95\% confidence interval for a population proportion is (0.67, 0.83).
	\begin{mcquestions}
		\question What is the sample proportion?
		\begin{mcanswerslist}
			\answer 0.67
			\answer 0.83
			\answer[correct] 0.75
			\answer 0.08
		\end{mcanswerslist}
		
		\question What is the margin of error?
		\begin{mcanswerslist}
			\answer 0.67
			\answer 0.83
			\answer 0.75
			\answer[correct] 0.08
		\end{mcanswerslist}
	\end{mcquestions}
	
	\begin{mcquestions}
		\question A 99\% confidence interval for a population proportion is $0.20 \pm 0.09$. What is the confidence interval?
		\begin{mcanswerslist}
			\answer (0.155,0.245)
			\answer (0.21,0.22)
			\answer[correct] (0.11,0.29)
			\answer The confidence interval cannot be determined.
		\end{mcanswerslist}
		
		\question The spinning dancer (or silhouette illusion) is a moving image of a woman that appears to be spinning.  Some people see her spinning clockwise and some see her spinning counterclockwise.  A student showed other students this and found that 30 out of 50 (or 60\%) saw her spinning clockwise.  The student researcher was interested in the probability of a randomly chosen student that would see the dancer spinning clockwise.  The student created the following null distribution.
		\includegraphics[width=\textwidth]{../images/null_dancer.jpg}
		Using the information provided above, determine a 2SD confidence interval for the long-run probability a randomly chosen student would see the woman spinning clockwise.  (Don't use any other applet to do this, just the information above.)
		\begin{mcanswerslist}
			\answer (0.36, 0.64)
			\answer[correct] (0.46, 0.74)
			\answer (0.53, 0.67)
			\answer (0.43, 0.57)
		\end{mcanswerslist}
		
		\question Many studies have investigated the question of whether people tend to think of an odd number when asked to think of a single-digit number (0 through 9; 0 is considered an even number). When asked to pick a number between 0 and 9, out of 70 students, 42 chose an odd number. Let the parameter of interest, $\pi$ , represent the probability that a student will choose an odd number.\\
		Are conditions met to use theory-based methods for conducting a 95\% confidence interval?
		\begin{mcanswerslist}
			\answer Yes, because there were two possible outcomes.
			\answer Yes, since the number of students is larger than 30.
			\answer[correct] Yes, since the number of students who chose an odd number, and the number of students who chose an even number, are both greater than 10.
			\answer No, since the sample size is too small.
		\end{mcanswerslist}
		
		\question Hill and Barton (Nature, 2005) conducted a study to investigate whether Olympic athletes in certain uniform colors have an advantage over their competitors. They noted that competitors in the combat sports of boxing, tae kwon do, Greco-Roman wrestling, and freestyle wrestling are randomly assigned to wear red or blue uniforms.  For each match in the 2004 Olympics, the uniform color of the winner was recorded. Hill and Barton found that in the 457 matches, the competitor wearing red won 248 times (54.3\%), whereas the person wearing blue won 209 times (45.7\%).\\
		Based on these data, use the 2SD method to construct a 95\% confidence interval for the proportion of matches won by competitors wearing red uniforms.
		\begin{mcanswerslist}
			\answer[correct] (0.497, 0.589)
			\answer (0.454, 0.546)
			\answer (0.493, 0.593)
			\answer (0.483, 0.602)
		\end{mcanswerslist}
		
		\question Suppose a 95\% confidence interval for a population proportion is (0.21, 0.43). How would the width of the confidence interval change if the confidence level was changed?
		\begin{mcanswerslist}
			\answer Both a 99\% confidence interval and a 90\% confidence interval would both be wider.
			\answer[correct] A 99\% confidence interval would be wider and a 90\% confidence interval would be narrower.
			\answer A 99\% confidence interval would be narrower and a 90\% confidence interval would be wider.
			\answer Both a 99\% confidence interval and a 90\% confidence interval would be narrower.
		\end{mcanswerslist}
		
	\end{mcquestions}
	
	Researchers are interested in investigating the effect of a drug used in the treatment of patients with glaucoma (an eye disorder associated with high eye pressure). In a volunteer sample of 35 patients with glaucoma in both eyes, one eye of each patient was randomly assigned to this drug, and the other eye was given a placebo. After one week, the eye pressure was measured on each eye. The difference in eye pressure between the two eyes (drug – placebo) was measured for each patient. The sample mean difference in eye pressure was –1.21 mmHg (millimeters of mercury), and the sample standard deviation of the differences was 4.67 mmHg.
	\begin{mcquestions}
		\question Use the 2SD method to find a 95\% confidence interval for the population mean difference in eye pressure (drug – placebo).
		\begin{mcanswerslist}
			\answer  (-2.51, 0.09)
			\answer  (-2.00, -0.42)
			\answer[correct]  (-2.79, 0.37)
			\answer  Not enough information exists.
		\end{mcanswerslist}
		
		\question Suppose a 95\% confidence interval for the population mean difference in eye pressure is (L, U). How would you interpret this interval?
		\begin{mcanswerslist}
			\answer  95\% confidence exists that each patient's difference in eye pressure (drug – placebo) will be between L and U mmHg.
			\answer[correct]  95\% confidence exists that the population mean difference in eye pressure (drug – placebo) if treated would be between L and U mmHg.
			\answer  Approximately 95\% of patients will have a difference in eye pressure (drug – placebo) between L and U mmHg.
			\answer  In the long run, 95\% of samples will produce a confidence interval of L to U mmH.
		\end{mcanswerslist}
		
		\question A 95\% confidence interval is computed to estimate the mean household income for a city. Which of the following values will definitely be within the limits of this confidence interval?
		\begin{mcanswerslist}
			\answer  The population mean
			\answer[correct]  The sample mean
			\answer  The sample standard deviation
			\answer  The p-value
		\end{mcanswerslist}
		
		\question On October 5, 2005, a tour boat named the Ethan Allen capsized on Lake George in New York with 47 passengers aboard. In the inquiries that followed, one suggestion was that the tour operators should have realized that the combined weight of so many passengers was likely to exceed the weight capacity of the boat, 7500 lbs.
		Based on this information, how surprising should the inquirers be for a sample of 47 passengers to have an average weight of at least 7500/47 = 159.57 lbs?
		\begin{mcanswerslist}
			\answer[correct]  Not surprising at all – the probability is over 90\%.
			\answer  Not surprising – the probability is between 50\% and 60\%.
			\answer  Somewhat surprising – the probability is between 20\% and 30\%.
			\answer  Very surprising – the probability is less than 5\%.
		\end{mcanswerslist}
	\end{mcquestions}
\end{document}